% figures_inventory_unified.tex
% Unified inventory of every figure embedded (via \includegraphics) anywhere
% in the submission: the main manuscript (jmlr_paper_main.tex) and the
% supplementary/erratum document set. Merges figures_inventory.tex and
% figures_inventory_supplementary.tex into a single source of record.

\begin{table}[H]
\centering
\caption{Figures embedded in the compiled main manuscript, in order of
  appearance. All 5 files are required in \texttt{submission/sources/figures/}.}
\label{tab:figlist_main}
\small
\renewcommand{\arraystretch}{1.25}
\begin{tabular}{L{5.2cm} L{3.0cm} L{5.3cm}}
\toprule
\textbf{Filename} & \textbf{Location} & \textbf{Content}\\
\midrule
\texttt{hypatiax\_three\_systems} &
  Methodology &
  HypatiaX architecture: LLM symbolic generation, residual MLP, and
  routing/ensembling layer.\\
\texttt{hypatiax\_algorithm1\_routing\_cascade\_v2} &
  Methodology, Component~3 &
  Algorithm~1 routing cascade diagram (trust/extrapolation detection,
  LLM query, PySR phases, convergence lock).\\
\texttt{fig18\_r2\_heatmap\_improved} &
  Results, Core-15 benchmark &
  $\Rsq$ heatmap (PySR-only vs.\ HypatiaX), values clipped to
  $[-1,1]$.\\
\texttt{fig09\_r2\_heatmap\_regimes} &
  Results, Core-15 benchmark &
  $\Rsq$ heatmap with raw (unclipped) values, showing severity of
  extrapolation failures.\\
\texttt{fig1\_seed\_sweep} &
  Results, Portfolio Variance &
  Seed sweep: Near/Medium/Far $\Rsq$ for PySR-only vs.\ HypatiaX
  across 5 seeds.\\
\bottomrule
\end{tabular}
\end{table}

\begin{table}[H]
\centering
\caption{Figures embedded across the supplementary/erratum document set. Only
  \textbf{1} distinct figure exists; it appears in two files.}
\label{tab:figlist_supplementary}
\small
\renewcommand{\arraystretch}{1.25}
\begin{tabular}{L{5.2cm} L{3.6cm} L{5.3cm}}
\toprule
\textbf{Filename} & \textbf{Appears in} & \textbf{Content}\\
\midrule
\texttt{recovery\_vs\_noise} &
  Standalone file; also embedded as Figure~1 (p.~6) of
  \texttt{jmlr\_erratum\_report.pdf} &
  Recovery rate ($R^2\ge0.999999$) vs.\ noise level $\sigma\in\{0,0.05,0.1,
  0.5,1\}\%$, July 2026 rerun, $n{=}200$, 30 equations, all 6 methods
  (PureLLM, HLLMNN, SymLLM, EHSDeFi, HDSv50\_2, ImpNN).\\
\bottomrule
\end{tabular}
\end{table}

\noindent\textbf{Combined total:} \textbf{6} distinct figure files across the
entire submission --- 5 in the main manuscript (Table~\ref{tab:figlist_main})
and 1 in the supplementary/erratum set (Table~\ref{tab:figlist_supplementary}).
The two sets are disjoint: none of the 5 main-manuscript figures recur in the
supplementary set, and \texttt{recovery\_vs\_noise} does not appear anywhere in
\texttt{jmlr\_paper\_main.tex}.

\noindent\textbf{Scope and verification:} the main-manuscript list is
exhaustive for the compiled PDF --- every \verb|\includegraphics| call in
\texttt{jmlr\_paper\_main.tex} is accounted for in
Table~\ref{tab:figlist_main}, and neither supplementary document embeds any
images beyond \texttt{recovery\_vs\_noise}. The supplementary list was checked
directly by rendering every page of all four supplementary/erratum PDFs
(\texttt{correction\_final\_report.pdf}, \texttt{hllmnn\_erratum.pdf},
\texttt{jmlr\_erratum\_report.pdf}, \texttt{recovery\_vs\_noise.pdf}) and
confirming with \texttt{pdfimages -list} that none contain embedded raster
images --- the one figure present is a vector plot.
\texttt{correction\_final\_report.pdf} and \texttt{hllmnn\_erratum.pdf}
contain tables only, zero figures. \texttt{table6\_comparison.tex} likewise
contains no \verb|\includegraphics| call --- it is a colour-coded comparison
table, not a figure source.

\noindent\textbf{Important --- avoid duplicate submission:}
\texttt{recovery\_vs\_noise.pdf} is \emph{not} a second, independent figure.
It is the same chart embedded as Figure~1 in \texttt{jmlr\_erratum\_report.pdf}
(identical title, axes, data, and six-method legend, confirmed by direct
visual comparison of both renders). If both the erratum PDF and its
constituent figure are submitted as supplementary material, this should be
stated explicitly so a reader does not count it as two figures, or reviewers
do not ask why a ``new'' figure appears unexplained outside the erratum text.

\noindent\textbf{Generation source:} unlike the benchmark supplement's
separate reproducibility figure set (regenerated from
\texttt{noise\_sweep}/\texttt{sample\_complexity} v2 JSONs via
\texttt{generate\_figures.py}), no generation script or source data file for
the 5 main-manuscript figures is currently documented in the supplementary
materials. If one exists, add it here; otherwise treat the image files
themselves as the source of record.
